% =====================================================================
%  Lecture 12 -- Mobility Management Entity (MME)         (English edition)
% =====================================================================
\documentclass[11pt,a4paper]{article}
\providecommand\ACNroot{../..}
\input{\ACNroot/templates/lecture-en}

\ACNset{number}{12}
\ACNset{title}{Mobility Management Entity (MME)}
\ACNset{shorttitle}{MME}
\ACNset{titleother}{موجودیت مدیریت تحرک (MME)}
\ACNset{part}{Part II -- 4G access and core (Lectures 5--13)}
\ACNset{time}{6--8 hours (3 h reading, 3--5 h exercises)}
\ACNset{level}{Intermediate}
\ACNset{prereq}{Lectures 4, 6, 9 and 11; basic queueing (M/M/1)}
\ACNset{gen}{4}

\begin{document}
\ACNcover

% ---------------------------------------------------------------------
\section{Learning objectives}
After studying this lecture you should be able to:
\begin{objectives}
  \item \textbf{List} the functions and reference points of the MME and \textbf{explain} why it is a pure
        control-plane node.
  \item \textbf{Decode} a GUTI and \textbf{explain} how temporary identities protect privacy and route
        signalling inside an MME pool.
  \item \textbf{Trace} the initial attach procedure end to end.
  \item \textbf{Evaluate} paging strategies by their expected signalling cost and delay.
  \item \textbf{Dimension} an MME pool with weight factors and failure margins, and \textbf{analyse} MME
        overload with a queueing model.
  \item \textbf{Compare} the MME with the 2G/3G SGSN/VLR and the 5G AMF.
\end{objectives}

% ---------------------------------------------------------------------
\section{Prerequisite knowledge}
You need the EPC architecture and EMM/ECM states of Lecture~9, the S1AP and RRC procedures of
Lecture~6, location-area design from Lecture~4 and the S6a dialogue with the HSS (Lectures~10--11).
The overload model uses the M/M/1 queue: Poisson arrivals at rate $\lambda$, exponential service at
rate $\mu$, one server.

\paragraph{Key terms} The vocabulary of this lecture with its Persian equivalents:
\begin{keyterms}
MME & \fa{موجودیت مدیریت تحرک} & Control-plane node of the EPC that terminates NAS signalling. \\
NAS & \fa{لایهٔ غیردسترسی} & Signalling between UE and MME (EMM and ESM) \cite{ts24301}. \\
GUTI & \fa{شناسهٔ موقت یکتای سراسری} & Temporary UE identity allocated by the MME. \\
S-TMSI & \fa{TMSI کوتاه‌شده} & MMEC + M-TMSI; used in paging and RRC set-up. \\
TAI list & \fa{فهرست شناسه‌های ناحیهٔ ردیابی} & Set of tracking areas in which the UE need not report. \\
Paging & \fa{فراخوانی} & Locating an idle UE for mobile-terminated traffic. \\
MME pool & \fa{مخزن MME} & Group of MMEs serving the same pool area (S1-flex). \\
Relative MME capacity & \fa{ظرفیت نسبی MME} & Weight used by eNBs to share new UEs among pool members. \\
Overload control & \fa{کنترل اضافه‌بار} & Mechanisms that reject or delay signalling when the MME is saturated. \\
T3412 & \fa{زمان‌سنج TAU دوره‌ای} & Periodic tracking-area update timer. \\
\end{keyterms}

% ---------------------------------------------------------------------
\section{Introduction and motivation}
In GPRS and UMTS the SGSN did everything for packet subscribers: it authenticated them, tracked them,
held their sessions and forwarded their packets. The EPS split this node into two. The
\emph{Mobility Management Entity} (MME) keeps the control functions; the serving gateway keeps the
packets \cite{ts23401}. The split lets operators scale signalling capacity (which grows with the number
of devices and their mobility) independently of user-plane capacity (which grows with traffic volume),
and it is the first step of the control/user-plane separation that 5G completes.

The MME is the UE's counterpart in the core: every NAS message -- attach, tracking-area update, service
request, bearer request -- terminates there. This lecture looks at what the MME stores, how it locates
idle devices, how MMEs share load in a pool, and what happens when signalling exceeds capacity.

% ---------------------------------------------------------------------
\section{Fundamental concepts}

\subsection{Functions of the MME}
\begin{itemize}
  \item \textbf{NAS signalling and security}: EMM and ESM procedures, authentication with vectors from
        the HSS, NAS ciphering and integrity \cite{ts24301,ts33401}.
  \item \textbf{Reachability and mobility}: tracking-area lists, paging of idle UEs, inter-eNB handover
        with MME relocation (S10), inter-RAT mobility to 2G/3G (S3).
  \item \textbf{Session management}: default and dedicated bearer set-up via the S-GW (S11), selection of
        S-GW and P-GW.
  \item \textbf{Subscription handling}: registration with the HSS (S6a), enforcement of access
        restrictions, equipment-identity check (S13).
  \item \textbf{Interworking}: CS fallback and SMS over SGs \cite{ts23272}, SRVCC over Sv \cite{ts23216},
        location services (SLs).
\end{itemize}
The MME never touches user packets; its load is measured in procedures per second and in stored
contexts, not in bits per second.

\begin{table}[htbp]
\centering\small
\caption{MME reference points.}\label{tab:ifaces}
\rowcolors{2}{acnmist}{white}
\begin{tabularx}{\linewidth}{@{}l l l Y@{}}
\toprule
\rowcolor{white}Interface & Peer & Protocol & Purpose \\
\midrule
S1-MME & eNB      & S1AP over SCTP & UE-associated and non-UE signalling, NAS transport \cite{ts36413} \\
S11    & S-GW     & GTPv2-C        & Session and bearer management \cite{ts29274} \\
S6a    & HSS      & Diameter       & Authentication data, location update \cite{ts29272} \\
S10    & MME      & GTPv2-C        & Context transfer on MME relocation \\
S3     & SGSN     & GTPv2-C        & Inter-RAT mobility with 2G/3G \\
S13    & EIR      & Diameter       & Equipment identity (IMEI) check \\
SGs    & MSC/VLR  & SGsAP over SCTP & CS fallback, SMS \cite{ts23272} \\
Sv     & MSC      & GTPv2-C        & SRVCC handover of voice to CS \cite{ts23216} \\
\bottomrule
\end{tabularx}
\end{table}

\begin{acnfigure}
  \acnfig{lecture-12/mme-interfaces}
  \caption{The MME and its reference points. Dashed lines are control-plane interfaces.}
  \label{fig:ifaces}
\end{acnfigure}

\subsection{Identities}
The permanent IMSI is sent over the air as rarely as possible. After the first attach the MME allocates
a \emph{Globally Unique Temporary Identity} (GUTI) \cite{ts23003} (\cref{fig:guti}). Its first part,
the GUMMEI, names the MME that allocated it (PLMN, 16-bit MME group ID, 8-bit MME code); its last part,
the 32-bit M-TMSI, names the UE within that MME. The 40-bit S-TMSI (MME code + M-TMSI) is enough inside
one pool and is used in paging and RRC connection set-up. Reallocating the GUTI regularly prevents an
observer from linking a UE's transactions.

\begin{acnfigure}
  \acnfig{lecture-12/guti}
  \caption{Structure of the GUTI and the identities derived from it.}
  \label{fig:guti}
\end{acnfigure}

\subsection{MME pools}
An eNB can connect to several MMEs (S1-flex). The MMEs that serve the same set of tracking areas form a
\emph{pool}. When a UE first connects, the eNB's NAS node selection function chooses an MME, weighted by
the \emph{relative MME capacity} each MME announced in S1 Setup; afterwards the MME code in the S-TMSI
routes every new connection of that UE to the same MME. Pools provide load sharing, resilience (a failed
MME's UEs re-register with the others) and the ability to take an MME out of service gracefully.

% ---------------------------------------------------------------------
\section{Detailed technical explanation}

\subsection{UE context}
For each registered UE the MME stores the IMSI and GUTI, the security context (keys, NAS counts,
algorithms), the subscription profile from the HSS, the TAI list and last known TAI and eNB, the EMM/ECM
state, the selected S-GW and the bearer contexts (EBI, QoS, S-GW TEIDs). In ECM-CONNECTED the context
also holds the S1AP identifiers of the UE's eNB connection. This state is what makes the MME the
\emph{anchor} for the UE's signalling: losing it (e.g.\ after an MME failure) forces the UE to attach
again.

\subsection{Tracking areas and paging}
The MME gives each UE a list of tracking areas (TAI list); the UE performs a TAU only when it enters a
TA outside the list or when the periodic timer T3412 expires (default 54\,min) \cite{ts24301}. When
downlink data arrive for an idle UE, the S-GW sends Downlink Data Notification and the MME pages the UE
in the cells of its TAI list with its S-TMSI. Operators reduce paging load with \emph{staged} strategies
-- first the last-seen eNB, then the last TA, then the whole list -- at the price of extra delay
(\cref{sec:paging}).

If no periodic TAU arrives, the MME starts the \emph{mobile reachable timer} (slightly longer than T3412);
when it expires, the MME stops paging the UE, and after a further implicit-detach timer it detaches it.

\subsection{Gateway selection}
At attach, the MME selects an S-GW that serves the UE's tracking area and a P-GW that serves the
requested APN, normally by DNS queries on names built from the TAI and the APN, which return candidate
gateways with priorities and weights. The P-GW must stay the same for the life of the PDN connection
(it owns the IP address); the S-GW may change with mobility (Lecture~13).

\subsection{Load balancing, rebalancing and overload control}
Under normal conditions the eNBs distribute new UEs in proportion to the relative capacities. To empty an
MME (for maintenance) or to rebalance, the MME releases UEs with the cause \emph{load-balancing TAU
required}: the UE performs a TAU, and because its S-TMSI is withheld, the eNB selects a new MME.
When an MME is overloaded it sends S1AP \msg{Overload Start} to its eNBs, which then reject new RRC
connections of selected types (e.g.\ all except emergency and mobile-terminated) until \msg{Overload Stop}
\cite{ts36413}. NAS back-off timers and access-class barring (Lecture~8) complement this at the UE.

% ---------------------------------------------------------------------
\section{Mathematical models}

\subsection{Staged paging}\label{sec:paging}
\paragraph{Model} The MME pages in stages $s = 1,2,\dots$; stage $s$ covers $n_s$ cells and finds the UE
with (unconditional) probability $p_s$, with $\sum_s p_s = 1$. Each stage takes one paging response
time $T_\mathrm{p}$. The expected number of cells paged and the expected delay are
\begin{equation}\label{eq:paging}
  \E[M] = \sum_{s} n_s\Bigl(1-\sum_{j<s}p_j\Bigr), \qquad
  \E[D] = T_\mathrm{p}\sum_{s}\Bigl(1-\sum_{j<s}p_j\Bigr).
\end{equation}
\paragraph{Example} Three stages: last eNB ($n_1 = 1$, $p_1 = 0.6$), last TA ($n_2 = 20$, $p_2 = 0.3$),
TAI list ($n_3 = 60$, $p_3 = 0.1$): $\E[M] = 1 + 0.4\times20 + 0.1\times60 = 15$ cells, compared with 60
for blanket paging; $\E[D] = 1.5\,T_\mathrm{p}$ instead of $T_\mathrm{p}$.
\paragraph{Interpretation} Staging pays off when the probabilities are concentrated in the early stages
(static users, small cells), and is unattractive for delay-sensitive services such as incoming voice
calls, for which operators often page the whole list at once.

\subsection{Pool weights and failure margin}\label{sec:pool}
\paragraph{Model} The share of new UEs that MME $i$ receives is
\begin{equation}\label{eq:weight}
  s_i = \frac{w_i}{\sum_j w_j},
\end{equation}
with $w_i$ its relative capacity. If a pool of $m$ equal MMEs must survive the loss of one without
overload, each MME may be loaded at most to
\begin{equation}\label{eq:pool}
  \rho_{\max} = \frac{m-1}{m}.
\end{equation}
\paragraph{Example} Weights $100, 100, 50$ give shares $0.4, 0.4, 0.2$. A pool of four equal MMEs may run
at 75\,\% in normal operation; a pool of two at only 50\,\%.

\subsection{Processing delay and overload}\label{sec:mm1}
\paragraph{Model} Treat the MME's signalling processor as an M/M/1 queue with capacity $\mu$ procedures per
second and offered load $\lambda$ ($\rho = \lambda/\mu < 1$). The mean response time is
\begin{equation}\label{eq:mm1}
  W = \frac{1}{\mu-\lambda} = \frac{1}{\mu(1-\rho)} .
\end{equation}
\paragraph{Example} $\mu = 5000$\,s$^{-1}$: $W = 1$\,ms at $\lambda = 4000$, 10\,ms at 4900 and 100\,ms at 4990
(\cref{fig:overload}).
\paragraph{Interpretation} The delay grows without bound as $\rho\to1$. Once response times approach NAS
timers (T3410 = 15\,s for attach), UEs retransmit, the offered load rises further and the system can
collapse; overload control must therefore act well before $\rho = 1$. The M/M/1 model is optimistic for
the bursty arrivals of a signalling storm.

\begin{acnfigure}
  \acnfig{lecture-12/overload}
  \caption{Mean response time of an M/M/1 signalling processor versus load (\cref{eq:mm1}).}
  \label{fig:overload}
\end{acnfigure}

% ---------------------------------------------------------------------
\section{Network architecture and diagrams}
In a typical network, the eNBs of a region connect to all MMEs of a pool of two to eight MMEs; pools of
neighbouring regions overlap so that a UE moving between them changes MME only at the pool boundary.
The MMEs reach the HSS via DRAs, the S-GWs and P-GWs over S11 (often in the same data centre), and the
MSC/VLR over SGs for CS fallback and SMS (\cref{fig:ifaces}).

% ---------------------------------------------------------------------
\section{Step-by-step procedures}

\begin{procedurebox}{Initial attach (simplified)}\label{proc:attach}
\begin{steps}
  \item The UE sends Attach Request with a PDN Connectivity Request inside RRC set-up; the eNB selects an
        MME and forwards it in S1AP Initial UE Message.
  \item The MME identifies the UE (GUTI or IMSI), runs EPS-AKA with a vector from the HSS and activates NAS
        security.
  \item The MME registers with the HSS (ULR/ULA) and receives the subscription profile.
  \item The MME selects S-GW and P-GW and sends Create Session Request; the P-GW allocates the IP address
        and, with a PCRF, obtains the policy; responses return the S-GW's uplink tunnel endpoint.
  \item The MME sends Initial Context Setup Request with Attach Accept and the default bearer to the eNB,
        which configures the radio bearer with RRC Connection Reconfiguration.
  \item The UE answers Attach Complete; the MME sends Modify Bearer Request with the eNB's downlink tunnel
        endpoint to the S-GW. The default bearer now carries traffic in both directions.
\end{steps}
\end{procedurebox}

\begin{acnfigure}
  \resizebox{\linewidth}{!}{\acnfig{lecture-12/attach}}
  \caption{Simplified E-UTRAN initial attach (\cref{proc:attach}) \cite{ts23401}.}
  \label{fig:attach}
\end{acnfigure}

\begin{procedurebox}{Emptying an MME for maintenance}\label{proc:rebalance}
\begin{steps}
  \item Set the MME's relative capacity to 0 (S1AP MME Configuration Update) so that eNBs stop selecting
        it for new UEs.
  \item Release connected UEs with cause \emph{load-balancing TAU required}; page idle UEs or wait for their
        next contact and release them in the same way.
  \item Each UE performs a TAU without S-TMSI routing; the eNB selects another MME, which fetches the
        context over S10 and registers with the HSS.
  \item When no contexts remain, the MME can be taken out of service.
\end{steps}
\end{procedurebox}

% ---------------------------------------------------------------------
\section{Worked examples}

\begin{examplebox}{Decoding a GUTI}\label{ex:guti}
\textbf{Problem.} A GUTI has MCC 432, MNC 35, MMEGI \texttt{0x8001}, MMEC \texttt{0x02} and M-TMSI
\texttt{0xC0A80001}. Which identities does the eNB use, and which MME serves the UE?

\textbf{Solution.} GUMMEI $=$ (432, 35, 32769, 2): MME code 2 in group 32769 of PLMN 432-35. The S-TMSI is
\texttt{0x02C0A80001} (40 bits); paging and RRC set-up carry only this. The eNB's NAS node selection
function maps MMEC 2 to the corresponding S1 connection in its pool.
\end{examplebox}

\begin{examplebox}{Paging cost}\label{ex:paging}
\textbf{Problem.} A city has $5\times10^5$ mobile-terminated events per busy hour. Compare the paging
messages per busy hour for blanket paging over 60 cells with the three-stage strategy of \cref{sec:paging}.

\textbf{Solution.} Blanket: $5\times10^5\times60 = 3.0\times10^7$ cell pagings per hour. Staged:
$5\times10^5\times15 = 7.5\times10^6$, a factor 4 less, at a mean delay of $1.5\,T_\mathrm{p}$.
\end{examplebox}

\begin{examplebox}{Pool sizing}\label{ex:pool}
\textbf{Problem.} A region has $6\times10^6$ registered UEs. Each MME supports $2\times10^6$ contexts. How many
MMEs are needed if the pool must survive one failure?

\textbf{Solution.} Without redundancy $\lceil 6/2\rceil = 3$. With one failure tolerated, the remaining $m-1$
MMEs must hold all contexts: $(m-1)\times2\times10^6 \ge 6\times10^6 \Rightarrow m = 4$, each loaded to
$6/8 = 75\,\%$ in normal operation, consistent with \cref{eq:pool}.
\end{examplebox}

% ---------------------------------------------------------------------
\section{Technology comparison}
\begin{table}[htbp]
\centering\small
\caption{Mobility-management nodes across generations.}\label{tab:cmp}
\rowcolors{2}{acnmist}{white}
\begin{tabularx}{\linewidth}{@{}l Y Y Y@{}}
\toprule
\rowcolor{white} & SGSN + MSC/VLR (2G/3G) & MME (4G) & AMF (+ SMF) (5G) \\
\midrule
User plane       & SGSN forwards packets (3G: optional direct tunnel) & None & None \\
Session control  & SGSN                        & MME (via S-GW)    & SMF (separate) \\
Location unit    & Location / routing area     & TAI list          & Registration area (TAI list) \\
Temporary ID     & TMSI / P-TMSI               & GUTI, S-TMSI      & 5G-GUTI, 5G-S-TMSI \\
Pooling          & Iu-flex (optional)          & S1-flex           & AMF sets, NRF \\
Interfaces       & MAP, Gn/Gp, Iu              & S1AP, GTPv2-C, Diameter & NGAP, HTTP/2 SBI \\
\bottomrule
\end{tabularx}
\end{table}

\paragraph{From MME to AMF} The 5G AMF keeps the MME's access and mobility functions but hands session
management to a separate SMF, and replaces GTPv2-C and Diameter with service-based interfaces
(Lecture~19). The concepts -- temporary identities, registration areas, pools, overload control -- carry
over almost unchanged.

% ---------------------------------------------------------------------
\section{Key takeaways}
\begin{takeawaybox}
\begin{itemize}
  \item The MME is the EPC's control-plane anchor: NAS, security, mobility, bearer control, gateway
        selection; no user data.
  \item GUTI = GUMMEI + M-TMSI; the S-TMSI routes signalling to the right MME in a pool and protects the IMSI.
  \item Attach combines authentication, HSS registration, session creation and radio-bearer set-up.
  \item Staged paging trades delay for signalling: $\E[M] = \sum n_s(1-\sum_{j<s}p_j)$.
  \item Pools share load by weights; $N+1$ redundancy limits normal load to $(m-1)/m$.
  \item Response time explodes as $\rho\to1$; overload control must act early.
\end{itemize}
\end{takeawaybox}

% ---------------------------------------------------------------------
\section{Exercises}

\begin{exercisebox}{1}{Control or user plane?}\label{exr:split}
For each function, state whether the MME, the S-GW or the P-GW performs it: (a) paging; (b) buffering
downlink packets for an idle UE; (c) IP address allocation; (d) NAS ciphering; (e) choosing the S-GW.
\end{exercisebox}

\begin{exercisebox}{1}{Why temporary identities?}\label{exr:guti}
Give two reasons why the network uses the GUTI/S-TMSI instead of the IMSI in paging and RRC set-up.
\end{exercisebox}

\begin{exercisebox}{2}{Paging strategy}\label{exr:paging}
Evaluate \cref{eq:paging} for $p = (0.4, 0.4, 0.2)$ and $n = (1, 30, 90)$. Compare with blanket paging over 90
cells. Which strategy would you use for incoming voice calls?
\end{exercisebox}

\begin{exercisebox}{2}{Weights after a failure}\label{exr:weights}
A pool has five MMEs with relative capacities $2,2,1,1,1$ (capacity proportional to weight). Compute the
shares of new UEs before and after the failure of one weight-2 MME, and the maximum normal-operation
utilisation that still allows the pool to absorb this failure.
\end{exercisebox}

\begin{exercisebox}{3}{Overload}\label{exr:overload}
An MME has $\mu = 5000$ procedures/s. (a) What load keeps $W \le 10$\,ms in the M/M/1 model? (b) After a
regional power outage, $2\times10^6$ UEs try to attach within 60\,s. Compare the offered load with $\mu$ and
explain which mechanisms keep the system stable.
\end{exercisebox}

\begin{exercisebox}{3}{Choosing T3412}\label{exr:t3412}
An MME holds $5\times10^6$ idle UEs. Compute the periodic-TAU rate for T3412 = 54\,min and 12\,min, and discuss
the trade-off between signalling, UE battery and the time needed to detect unreachable UEs. Why do IoT
devices use much longer values?
\end{exercisebox}

\subsection*{Solutions}
\addcontentsline{toc}{subsection}{Solutions}

\begin{solutionbox}{\ref{exr:split}}
(a) MME (triggered by the S-GW's Downlink Data Notification); (b) S-GW; (c) P-GW; (d) MME; (e) MME.
\end{solutionbox}

\begin{solutionbox}{\ref{exr:guti}}
Privacy: the IMSI is permanent, so sending it in the clear would let an eavesdropper track a subscriber;
the GUTI is reallocated. Efficiency and routing: the 40-bit S-TMSI is shorter than the IMSI and its MME code
tells the eNB which pool member holds the context.
\end{solutionbox}

\begin{solutionbox}{\ref{exr:paging}}
$\E[M] = 1 + 0.6\times30 + 0.2\times90 = 37$ cells versus 90; $\E[D] = (1 + 0.6 + 0.2)\,T_\mathrm{p} = 1.8\,T_\mathrm{p}$. For
voice, call set-up delay is critical and the UE must be found at once: page the whole TAI list (or at least
skip the single-cell stage).
\end{solutionbox}

\begin{solutionbox}{\ref{exr:weights}}
Before: $2/7 = 28.6\,\%$ for each weight-2 MME and $1/7 = 14.3\,\%$ for each weight-1 MME. After: total weight 5,
so $40\,\%$ for the remaining weight-2 MME and $20\,\%$ for each weight-1 MME. Capacity falls from 7 to 5 units,
so the normal utilisation must not exceed $5/7 = 71\,\%$.
\end{solutionbox}

\begin{solutionbox}{\ref{exr:overload}}
(a) $1/(\mu-\lambda) \le 0.01 \Rightarrow \lambda \le 4900$\,s$^{-1}$ ($\rho \le 0.98$); with bursty traffic the real
limit is lower. (b) $2\times10^6/60 = 3.3\times10^4$ attaches/s, almost seven times $\mu$ (and each attach involves
several messages). Stability requires S1AP overload control (eNBs reject or delay connections), access-class
barring and RACH back-off at the UE, randomised NAS retry timers, and a pool so that several MMEs share the
load.
\end{solutionbox}

\begin{solutionbox}{\ref{exr:t3412}}
$5\times10^6/3240 = 1543$\,TAU/s for 54\,min; $5\times10^6/720 = 6944$\,TAU/s for 12\,min. A short timer detects lost UEs
sooner (detection takes about T3412 plus the reachability margin) but multiplies signalling and drains batteries.
IoT devices send rarely and must last years on a battery; extended timers of many hours, combined with power
saving mode, reduce their signalling and energy use, and the network accepts that they are unreachable most of
the time.
\end{solutionbox}

% ---------------------------------------------------------------------
\section{Further reading}
The MME's role and procedures are specified in TS~23.401 \cite{ts23401}; NAS in TS~24.301 \cite{ts24301}; S1AP in
TS~36.413 \cite{ts36413}; GTPv2-C in TS~29.274 \cite{ts29274}; identities in TS~23.003 \cite{ts23003}. CS fallback
and SRVCC are in TS~23.272 \cite{ts23272} and TS~23.216 \cite{ts23216}. Olsson et al.\ \cite{olsson2013epc} and Sesia et
al.\ \cite{sesia2011lte} give textbook treatments; Kleinrock \cite{kleinrock1975queueing} covers the M/M/1 queue.

\printbibliography[heading=subbibliography,title={\ACNstrRefs}]
\end{document}
